Глава заполняется по мере появления runner'а сравнения моделей.
Ниже зафиксирована постановка эксперимента (сентябрь 2026~г.).

\section{Постановка}

Сравниваются детектор-бейзлайн (z-score / IQR по ряду) и три алгоритма без учителя: Isolation Forest, Local Outlier Factor, One-Class SVM.
Обучение не использует колонку \texttt{label}; метки NAB и KPI нужны только на этапе оценки.
Сплит~--- по времени внутри каждого ряда; для KPI тестовая часть берётся из \texttt{phase2\_ground\_truth}.

Метрики: PR-AUC, precision, recall (и при необходимости window-aware оценка по окнам NAB).
Prometheus в эту таблицу не входит: у точек нет эталона.

Ансамбль строится только если по ошибкам модели дополняют друг друга; иначе в продукт идёт одна лучшая модель.

\section{План прогонов}

\begin{enumerate}
  \item Каркас модуля \texttt{ml/} и тесты на небольших фикстурах.
  \item NAB: таблица моделей против бейзлайна.
  \item KPI: та же таблица на отложенном тесте.
  \item Качественные графики на рядах Prometheus (без PR-AUC).
\end{enumerate}

Результаты (таблицы, гиперпараметры, файлы моделей) будут добавлены в эту главу и в приложение.
